\documentclass{article}
\usepackage{mdposter}
\posttitle{Periods and rounding in integer optimization}
\postfield{Mathematics, number theory and optimization}

\begin{document}
\begin{mdposter}

\posterbox{name=c1, column=1, between=top and bottom}{%
\begin{mdcol}
\mdsec{The question}
Minimize a homogeneous polynomial $f$ of degree $d \ge 2$ with rational coefficients over the nonnegative
integer vectors $z$ with $Az = nb$, and call the minimum $m_f(n)$. If $v$ minimizes $f$ over the real relaxation, the
real minimum at scale $n$ is exactly $n^d f(v)$. We ask when $m_f(n)$ eventually agrees with a polynomial in $n$ on
each residue class, what the correction looks like when it does not, and how often rounding the real optimum to a
nearby integer point is exact.

\vspace{0.16in}
The same question arises in experimental design, where a design has to be realised in whole observations and
the natural candidate is the real optimum rounded off. Whether that rounding is right, and how the answer
depends on the scale, is the practical form of the question.
\mdgap
\mdsec{How we are approaching it}
\begin{itemize}
\item We show that integer optima stay within bounded distance of $nv$ when $f$ grows quadratically away from $v$.
\item We read the offsets as a point rotating on a torus, so the leading correction is a nearest-lattice function of that rotation.
\item We use the Subspace Theorem to bound how often an algebraic rotation lands near a wall where the answer switches.
\item We solve explicit families exactly, among them two-generator cubics through Pell equations.
\item We extend the classification to minima that sit on a face, where the period comes from the lattice of that face.
\item We verify the explicit families by exhaustive search at small scales before trusting the general argument.
\end{itemize}
\end{mdcol}
}

\posterbox{name=c2, column=2, between=top and bottom}{%
\begin{mdcol}
\mdsec{Nine classes, nine lines}
\includegraphics[width=\linewidth]{rigidity-classes.pdf}
\figcap{Figure 1.}{What the classification says, in one case. Minimizing $x^3+y^3$ under $x+4y=n$ has a rational real optimum of period 9, and the minimum then sits exactly on a straight line in $n$ within each residue class.}
\mdgap
\mdsec{What we have so far}
\begin{progress}
Suppose the feasible set is compact, $f$ has a unique minimizer $v$ with every coordinate positive, and the
Hessian at $v$ is positive definite on the kernel of $A$. Then $m_f$ is eventually quasipolynomial exactly when $v$
is rational, and its least period is the smallest $M$ with $Mv$ integral.
\end{progress}
\begin{progress}
Under the same hypotheses the scaled corrections converge. The limit takes finitely many values when $v$ is
rational and is a continuous law when $v$ is irrational.
\end{progress}
\begin{progress}
Choosing by the original cost among the minimizers of the quadratic term is exact for all $n \le N$ with at
most $O_\varepsilon(N^\varepsilon)$ exceptions, and at most $O(\log N)$ of them when the coordinates of $v$ generate
a quadratic field.
\end{progress}
\end{mdcol}
}

\posterbox{name=c3, column=3, between=top and bottom}{%
\begin{mdcol}
\mdsec{When rounding fails}
\includegraphics[width=\linewidth]{rigidity-rounding.pdf}
\figcap{Figure 2.}{Rounding the real optimum fails rarely, and the failures are exactly the terms of two recurrences. For $x^3+y^3$ under $2x+3y=n$ we found every failure up to $n = 20{,}000$ by exhaustive search. Each one costs exactly 1, and the count follows the predicted logarithmic rate.}
\mdgap
\mdsec{Open questions}
\begin{itemize}
\item We assume one real minimizer. What happens when the minimum is attained at several points?
\item Curvature cannot be dropped. We give polynomials with a unique rational minimizer whose integer minima are not eventually quasipolynomial, and we have no classification there.
\item The constants are ineffective, and we have no polynomial-time algorithm for the version on a face.
\item We count the exceptions exactly only for two generators. What is the sharp count in higher dimensions?
\end{itemize}
\end{mdcol}
}

\end{mdposter}
\end{document}
